\chapter*{Preface}
\addcontentsline{toc}{chapter}{Preface}
\markboth{Preface}{Preface}

\vspace{1em}

\noindent
A chiral algebra is a solution. What is the problem it solves?

Volumes~I and~II develop the bar-cobar machine that extracts
modular invariants from any chiral algebra~$A$ on a curve~$X$:
the bar complex $B(A) = T^c(s^{-1}\bar{A})$ with its
deconcatenation coproduct, the universal Maurer--Cartan element
$\Theta_A$, the modular characteristic
$\kappa_{\mathrm{ch}}(A)$, and the five theorems that control
the genus tower. Volume~I introduces $E_1$-chiral quantum groups
as abstract algebraic-geometric objects---ordered bar complexes
with spectral $R$-matrices, Drinfeld coproducts, and
Yang--Baxter equations, all living on configuration spaces of
curves. Neither volume constructs a single chiral algebra or
a single chiral quantum group. The input is always assumed.
This volume constructs the input: both the chiral algebras and,
from Calabi--Yau geometry, the concrete $E_1$-chiral quantum
groups that are the flesh on Volume~I's algebraic skeleton.

A terminological point: ``Hochschild'' in this volume means
categorical Hochschild cohomology: the cyclic bar complex
$\mathrm{CC}_*(\cC)$ of a dg or $\Ainf$-category, topological
in nature. The chiral upgrade $\mathrm{ChirHoch}^*(A)$ of
Volume~I (Theorem~H) is the holomorphic version; it is the
output of the functor $\Phi$, not the input. The two are
related by the functor itself: $\Phi$ sends $\mathrm{CC}_*(\cC)$
to $\mathrm{ChirHoch}^*(A_\cC)$, lifting topological Hochschild
data to holomorphic chiral data on the curve.

The construction is a functor
\[
 \Phi \colon \CY_d\text{-}\Cat \longrightarrow
 \Etwo\text{-}\mathrm{ChirAlg}\,.
\]
A Calabi--Yau category $\cC$ of dimension~$d$ carries a cyclic
$\Ainf$-structure: a non-degenerate trace
$\Tr \colon \HC^{-}_{d}(\cC) \to k$ on negative cyclic
homology, encoding the CY condition as a $d$-dimensional
Frobenius structure at the chain level. The functor $\Phi$
converts this datum into a chiral algebra $A_\cC$ on~$X$ in
three steps: the cyclic $\Ainf$-structure determines a Lie
conformal algebra~$L_\cC$; the factorization envelope
$U_X^{\mathrm{fact}}(L_\cC)$ produces a chiral algebra
$A_\cC$; the $\Sd$-framing enhances the factorization
structure. For $d = 2$, the $\mathbb{S}^2$-action on
Hochschild homology gives an $\Etwo$-chiral algebra with
braided monoidal representation category, and all three steps
are proved (Theorem~CY-A). For $d = 3$, holomorphic
Chern--Simons breaks $\Etwo$ to~$\Eone$; the braiding is
recovered via the Drinfeld center of the $\Eone$-monoidal
representation category. Steps~(1) and~(2) are unconditional;
step~(3) is a programme conditional on the chain-level
$\mathbb{S}^3$-framing.

\bigskip

\noindent\textbf{What the CY perspective reveals.}
Volumes~I and~II see one chiral algebra at a time. The CY
functor sees a category $\cC$ that admits \emph{multiple}
chiral algebraizations, each with its own modular
characteristic. A single CY$_3$ manifold~$X$ produces a
spectrum of $\kappa_\bullet$-values, not a single number. For
$K3 \times E$:
\[
 \operatorname{Spec}_{\kappa_\bullet}(K3 \times E)
 = \{\kappa_{\mathrm{cat}},\, \kappa_{\mathrm{ch}},\,
 \kappa_{\mathrm{BKM}},\, \kappa_{\mathrm{fiber}}\}
 = \{2,\, 3,\, 5,\, 24\}\,.
\]
Here $\kappa_{\mathrm{cat}} = 2 = \chi(\cO_{K3})$ is the
holomorphic Euler characteristic;
$\kappa_{\mathrm{ch}} = 3 = \dim_\C X$ is the modular
characteristic of the chiral de~Rham complex;
$\kappa_{\mathrm{BKM}} = 5$ is the weight of the Igusa cusp
form $\Delta_5$ (so $2\kappa_{\mathrm{BKM}} = 10 = \mathrm{wt}(\Phi_{10})$,
the Siegel modular weight); and
$\kappa_{\mathrm{fiber}} = 24$ is the lattice rank. These
four numbers measure four faces of the same manifold: the
holomorphic Euler characteristic, the boundary chiral algebra,
the bulk Borcherds--Kac--Moody superalgebra, and the lattice
frame. The chiral-holographic pair
$(\kappa_{\mathrm{ch}}, \kappa_{\mathrm{BKM}}) = (3, 5)$ is
the one that sees the boundary--bulk axis of Volumes~I--II.
The $\kappa_\bullet$-spectrum is invisible from any single
chiral algebra; it is intrinsic to the CY category.

The $\En$-hierarchy provides the organizing spiral. A CY$_d$
category carries a native $\Sd$-framing, hence an
$E_d$-chiral algebra via $\Phi$. At $d = 3$, the chiral
algebra is natively $E_3$; but
$\pi_1(\Conf_2(\R^3)) = 0$, so the $E_3$-braiding is
symmetric; too symmetric for quantum groups. The resolution
is a two-step detour: stabilize from $E_3$ to $\Eone$ (forget
the symmetric braiding, retain the associative product), then
recover a \emph{non-symmetric} braiding via the Drinfeld
center $Z(\Rep^{\Eone}(A_\cC))$, which is $\Etwo$-monoidal.
The resulting $R$-matrix is the quantum group $R$-matrix. The
key arithmetic: $\pi_1(\Conf_2(\R^2)) = \Z$ (braid group,
native $\Etwo$, non-symmetric), while
$\pi_1(\Conf_2(\R^d)) = \Z/2$ for $d \geq 3$ (symmetric
group). Dimension $d = 2$ is the sweet spot where the native
$\Etwo$-structure already carries a non-symmetric braiding.
For $d \geq 3$, the descent-and-center detour is mandatory:
one must discard the over-symmetric $E_d$-structure, pass to
$\Eone$, and rebuild the braiding categorically.

The spiral, not circle, because the closing step is
conjectural at $d \geq 3$. The trajectory
$E_d \to \Eone \to \Etwo$ (via the Drinfeld center) returns
to the same operadic level as $d = 2$, but the $\Etwo$
structure so obtained is \emph{derived}: it lives on the
representation category, not on the algebra itself. Whether
this derived $\Etwo$-structure lifts to a genuine $\Etwo$-chiral
algebra (closing the spiral to a circle) is the content of
Conjecture~CY-B. At $d = 2$ the spiral closes by construction:
the native $\Etwo$-framing and the Drinfeld center $\Etwo$
agree. At $d = 3$ the two $\Etwo$-structures inhabit different
categories, and identifying them requires the chain-level
$\mathbb{S}^3$-framing that is the open core of the $d = 3$
programme. The spiral is the honest shape of the $\En$-story:
each CY dimension $d$ enters at $E_d$, descends to $\Eone$,
and ascends to $\Etwo$ through the center; the question is
whether the ascent reaches the starting altitude.

Three further structures emerge only through $\Phi$. First,
the bar complex converts enumerative geometry into homotopy
algebra: Donaldson--Thomas invariants become root
multiplicities of the Borcherds--Kac--Moody superalgebra, and
the denominator identity becomes the bar Euler product.
Second, the $\Eone$-stabilization theorem: for $d \geq 3$,
the CY chiral algebra is natively $\Eone$ (the
Schouten--Nijenhuis bracket vanishes, all noncommutative
structure arises from quantization, and the CoHA extension
correspondence imposes a preferred ordering). The braided
monoidal structure lives not in the algebra itself but in
the Drinfeld center of its $\Eone$-representation category;
the Drinfeld center is the categorified averaging map
$\mathrm{av} \colon \frakg^{\Eone} \to \frakg^{\mathrm{mod}}$
of Volume~I. Third, the framing obstruction tower: the
effective obstruction
$\Obs_{\mathrm{eff}}(d) \in \pi_d(BU)$ is $8$-periodic by
Bott periodicity and trivial precisely when
$d \bmod 8 \in \{1, 3, 7\}$, giving a rigid topological
constraint on which CY dimensions admit additional chiral
structure beyond~$\Eone$.

\bigskip

\noindent\textbf{The prototype: $K3 \times E$.}
The product of a K3 surface~$S$ and an elliptic curve~$E$ is
the canonical CY$_3$ testing ground: simple enough to compute
(lattice $\Lambda^{3,2}$ of signature $(3,2)$, K3 elliptic
genus $\phi_{0,1}(\tau,z)$, Borcherds denominator
$\Delta_5$), rich enough to obstruct. Four structural
barriers separate the shadow obstruction tower from the full
automorphic form:
\begin{enumerate}[label=\textup{(O\arabic*)},nosep]
\item \emph{Categorical}: the shadow tower produces numbers
 (Taylor coefficients of $\Theta_{A_\cC}$); the denominator
 identity is a function (an automorphic form on $\cH_2$).
\item \emph{$\kappa_{\mathrm{ch}}$--$\kappa_{\mathrm{BKM}}$
 mismatch}: $\kappa_{\mathrm{ch}} = 3 \neq 5 =
 \kappa_{\mathrm{BKM}}$; the single-copy chiral algebra sees
 $\dim_\C X$, while the bulk reconstruction sees the
 Borcherds weight.
\item \emph{Second quantization}: the physical DT partition
 function is the DMVV symmetric product of single-copy data;
 the bar complex of a single algebra is the single copy.
\item \emph{Schottky}: at $g \geq 4$, the Torelli map
 $\cM_g \to \cA_g$ is no longer dominant, and the shadow
 tower is imprisoned in tautological classes on
 $\overline{\cM}_g$.
\end{enumerate}
Each obstruction points to a research programme. The
$\kappa_{\mathrm{ch}}$--$\kappa_{\mathrm{BKM}}$ mismatch~(O2)
is the most telling: a single CY$_3$ admits multiple chiral
algebraizations, and the $\kappa_\bullet$-spectrum records
which face of the geometry each algebraization captures.

The functor $\Phi$ applied to a K3 surface distinguishes
three algebras: the lattice vertex algebra $V_\Lambda$
attached to $H^*(K3, \Z)$, the generalized Kac--Moody algebra
whose denominator is $\Delta_5$, and the Conway module arising
from the Leech lattice frame. These three algebraizations
share $\kappa_{\mathrm{cat}} = 2$ but differ in
$\kappa_{\mathrm{ch}}$ and $\kappa_{\mathrm{BKM}}$; the
distinction is intrinsic to the CY functor and invisible at
the level of the Euler characteristic alone.

Thirteen structural results (K3-1 through K3-13) are proved for
$K3 \times E$; here are the five that most illuminate the
programme:
\emph{K3-1}: $\kappa_{\mathrm{ch}}(K3 \times E) = 3$, verified
by six independent paths (direct Hodge, Mukai pairing,
Dolbeault contracting homotopy, bar computation, Borcherds
lift, and elliptic genus).
\emph{K3-4}: the moonshine multiplier
$A_n = \kappa_{\mathrm{cat}} \cdot \dim(\rho_n)$, factoring BPS
degeneracies into a homological part ($\kappa_{\mathrm{cat}} = 2$)
and a group-theoretic part ($\rho_n$, the Mathieu representation).
\emph{K3-7}: the shadow--Siegel gap theorem with four structural
obstructions (O1--O4), proving that the shadow tower does NOT
produce $\Phi_{10}$.
\emph{K3-9}: the denominator identity
$\prod_{\alpha > 0}(1 - e^{-\alpha})^{\mathrm{mult}(\alpha)}$
equals the bar Euler product, conditional on CY-A for $d = 3$.
\emph{K3-12}: the K3 double current algebra $\fg_{K3}$, the
finite-dimensional analogue of the DDCA in which
$\C[u,v]$ is replaced by $H^*(S,\C)$ with the Mukai pairing.

\bigskip

\noindent\textbf{The holomorphic Chern--Simons hierarchy.}
The theories studied in this volume form a hierarchy in real
dimension, each producing chiral algebras by restriction to
codimension-$2$ defects. At $d = 6$: holomorphic Chern--Simons
on a Calabi--Yau $3$-fold~$Y$ (Witten, Costello), the primary
source. A curve $C \subset Y$ of codimension~$2$ carries a defect
algebra: the normal bundle $N_{C/Y}$ of rank~$2$ supplies two
spectral parameters, and the defect algebra is a chiral algebra
on~$C$ with quantum group structure from the normal directions
(Proposition~\ref{prop:codim2-defect-ope}). For $Y = \C^3$ and
$C = \C$, the defect algebra is $\cW_{1+\infty}$ at
$\Psi = -\sigma_2(h_1, h_2, h_3)$ (the second elementary
symmetric polynomial in the $\Omega$-background parameters).
At $d = 5$: partially topological Chern--Simons on $\R \times Y$
for a CY$_2$ surface~$Y$; restriction to $\R \times C$ gives the
ordered bar with one spectral parameter. At $d = 4$: holomorphic
Chern--Simons on $\C^2$; the Costello--Witten--Yamazaki theory
that produces Yangians from holomorphic gauge fields. At $d = 3$:
the theories of Volume~II. The hierarchy
$6 \to 5 \to 4 \to 3$ is codimension-$2$ restriction at each
step; each restriction produces a chiral algebra from a gauge
theory, and the $E_n$ level decreases as the dimension drops.
The hierarchy is not limited to Chern--Simons: 6d little string
theory on K3 and M5 brane theory produce chiral algebras by the
same defect mechanism, with the specific gauge-theory content
replaced by stringy data.

\bigskip

\noindent\textbf{Concrete CY quantum groups.}
Volume~I introduces $E_1$-chiral quantum groups as abstract
algebraic-geometric objects: an ordered bar complex
$B^{\mathrm{ord}}(A)$ with spectral $R$-matrix, Drinfeld
coproduct, and Yang--Baxter equation, all living on
configuration spaces of curves. This volume constructs the
concrete examples. The critical cohomological Hall algebra
$\CoHA(Q)$ of a quiver~$Q$ (with potential~$W$ from the CY
condition) is an $E_1$-chiral quantum group: its product is the
ordered Hall multiplication (graded by dimension vectors), its
$R$-matrix comes from the Ext-quiver, and its coproduct is the
CoHA vertex coproduct of Schiffmann--Vasserot. For the Jordan
quiver ($Q = \bullet \circlearrowleft$):
$\CoHA(Q) \cong Y^+(\widehat{\fgl}_1) \cong \cW_{1+\infty}$.
The Drinfeld center
$Z(\Rep^{E_1}(Y^+(\widehat{\fgl}_1))) \cong
\Rep^{E_2}(Y(\widehat{\fgl}_1))$ recovers the full braided
structure (Theorem~\ref{thm:c3-drinfeld-center}). For a general
toric CY$_3$ with fan~$\Sigma$: quiver charts indexed by
Bridgeland stability chambers give local chiral algebras;
wall-crossing is conjecturally MC gauge equivalence. The quantum
toroidal algebra $U_{q,t}(\widehat{\widehat{\fgl}}_1)$ with CY
condition $q_1 q_2 q_3 = 1$ is the two-parameter deformation
carrying the full $\Z \times \Z$ bigrading from the double
current algebra.

\bigskip

\noindent\textbf{Toric CY$_3$ and critical CoHAs.}
For a toric CY$_3$ determined by a fan~$\Sigma$, the toric
diagram fixes a quiver $Q_X$. The critical cohomological
Hall algebra $\CoHA(Q_X)$ is the $\Eone$-sector: the positive
half of an affine super Yangian
$Y(\widehat{\frakg}_{Q_X})$. The full braided structure is
recovered through the Drinfeld center. For $\C^3$: the
Jordan quiver gives
$Y(\widehat{\fgl}_1) \simeq \cW_{1+\infty}$, and the chain
$\C^3 \to \cW_{1+\infty} \to
\Rep^{\Etwo}(Y(\widehat{\fgl}_1))$ is verified end-to-end.
The toric fan is the Dynkin data.

The CoHA is the natural $\Eone$-primitive: it is the ordered
bar complex of the CY$_3$ chiral algebra, with deconcatenation
coproduct encoding the topological direction and the R-matrix
encoding the spectral braiding. The passage from CoHA to
quantum group is the CY incarnation of the
$\Eone$-to-$\Etwo$ lift: averaging over the symmetric group
recovers the factorization coalgebra of Volume~I.
The algebraic underpinning of this passage (the integrable
structure of the ordered chiral bar complex, the R-matrix, and
the Yangian presentation) is developed in the companion paper
\emph{On the Integrable Theory of Chiral Yangians and Related
Chiral Quantum Groups} (Volume~I, standalone).

\bigskip

\noindent\textbf{The shadow tower as Gromov--Witten theory.}
At $\kappa_{\mathrm{ch}} = \Psi$ (the deformation parameter of
$\cW_{1+\infty}$ at $c = 1$), the shadow obstruction tower of
Volume~I produces the perturbative constant-map Gromov--Witten
free energies $F_g^{\mathrm{GW,const}}(\C^3)$: the shadow IS
the full GW answer for $\C^3$ (which has no compact curves, hence
no worldsheet instantons). On the Donaldson--Thomas side, the
MacMahon function $M(q) = \prod_{n \ge 1} (1 - q^n)^{-n}$
encodes the DT partition function via the MNOP correspondence.
The shadow coefficients $S_r$ of the Virasoro algebra at the
GW specialisation are the $A_\infty$ coproduct correction
coefficients $\delta^{(r)}$: the shadow tower encodes the
coproduct corrections of the chiral quantum group.

The universal coproduct $\Delta_z(e_s) = \sum C(N_R{-}b, k)\,
z^k\, e_a^L \cdot e_b^R$ (all spins, closed form) extends the
Drinfeld coproduct from spin~$2$ to arbitrary spin, with the
universal Miura coefficient $(\Psi{-}1)/\Psi$ on the leading
cross-term $J \otimes W_{s-1}$ at every spin $s \ge 2$
(Theorem~\ref{thm:miura-cross-universality}).

\bigskip

\noindent\textbf{The $E_1$-chiral bialgebra.}
Volume~I defines $E_1$-chiral quantum groups abstractly. This
volume constructs them concretely from CY geometry. The
$E_1$-chiral bialgebra axioms (formalized in
\S\ref{sec:e1-chiral-bialgebra}) require: (H1)~an $E_1$-chiral
algebra $A$ on a curve $X$; (H2)~a spectral coproduct
$\Delta_z \colon A \to A \otimes A$ (living on the OPEN/$E_1$
colour); (H3)~an $R$-matrix $R(z)$ satisfying YBE; (H4)~RTT
compatibility; (H5)~a modular MC tower extending to all genera.
The ordered bar $B^{\mathrm{ord}}(A)$ preserves the $R$-matrix;
the symmetric bar $B^\Sigma(A)$ of Volume~I kills the Hopf
structure via averaging $\mathrm{av}(r(z)) = \kappa_{\mathrm{ch}}$.

The tetrahedron equation (the 3d analogue of YBE) fails for the
Yang $R$-matrix at $O(\kappa_{\mathrm{ch}}^2)$: the $E_3$
structure is genuinely nontrivial beyond $E_2$, confirming that
the $E_n$ operadic circle is not a formal consequence of
lower-level data.

\bigskip

\noindent\textbf{The Drinfeld centre as categorified averaging.}
Volume~I constructs the averaging map
$\mathrm{av} \colon \frakg^{E_1}_\cA \to
\frakg^{\mathrm{mod}}_\cA$ that projects the ordered
$R$-matrix to the scalar $\kappa_{\mathrm{ch}}$. This volume
constructs its categorical lift: the Drinfeld centre
$Z \colon E_1\text{-}\Cat \to E_2\text{-}\Cat$, which takes the
$E_1$-monoidal representation category
$\Rep^{E_1}(A_\cC)$ to an $E_2$-braided monoidal category
$Z(\Rep^{E_1}(A_\cC))$. Direct restriction from $E_3$ to $E_2$
gives only symmetric braiding
($\pi_1(\Conf_2(\R^3)) = 0$); the Drinfeld centre is the sole
source of non-trivial braiding. For $\C^3$:
$Z(\Rep^{E_1}(Y^+(\widehat{\fgl}_1)))
\cong \Rep^{E_2}(Y(\widehat{\fgl}_1))
\cong \Rep^{E_2}(\cW_{1+\infty})$
(Theorem~\ref{thm:c3-drinfeld-center}). Five invariants match
at six levels: naive centre (dimension~$1$) versus derived centre
(dimension~$3$, capturing $\Ext^{1,2}$ invisible to the
commutant). The Drinfeld centre is the deepest
construction of the programme: it closes the $E_n$ circle
by recovering $E_2$-braiding from $E_1$-ordering.

\medskip
\noindent\textbf{Wall-crossing as Maurer--Cartan gauge
equivalence.}\enspace
For a toric CY$_3$ with fan~$\Sigma$, Bridgeland stability
conditions carve the space of central charges into chambers.
Each chamber $\sigma$ determines a quiver presentation
$Q_\sigma$ and a critical CoHA
$\CoHA(Q_\sigma, W_\sigma)$. Wall-crossing (varying the
stability condition across a wall) changes the quiver
presentation but not the underlying CY category. The
conjecture: wall-crossing between chambers $\sigma$ and
$\sigma'$ is Maurer--Cartan gauge equivalence
$\Theta_\sigma \sim \Theta_{\sigma'}$ in the modular
convolution algebra of the CY chiral algebra. The global
chiral algebra $A_\cC$ would then be the homotopy colimit over
the chamber decomposition:
$A_\cC = \operatorname{hocolim}_\sigma A_{Q_\sigma}$.
For $\C^3$ (single chamber): the conjecture is vacuous. For
the resolved conifold $\cO(-1) \oplus \cO(-1) \to \P^1$
(two chambers): the wall-crossing formula is the flop relation
on DT partition functions (Bryan--Steinberg), and the MC
gauge equivalence is the bar-cobar transport across the flop.
This conjecture connects the algebraic programme of Volume~I
(Maurer--Cartan elements) to the enumerative programme
(DT/BPS wall-crossing) of this volume.

\bigskip

\noindent\textbf{From shadow tower to BPS degeneracies.}
The shadow obstruction tower of Volume~I produces scalar
invariants $S_r(\cA)$ at each degree~$r$. For CY$_3$ chiral
algebras, these invariants admit an enumerative interpretation:
the DT invariants $\Omega(\gamma)$ of the CY$_3$ manifold are
root multiplicities of the BKM superalgebra $\frakg_X$, and the
denominator identity
\[
\prod_{\alpha > 0} (1 - e^{-\alpha})^{\mathrm{mult}(\alpha)}
\;=\;
\sum_{w \in W} (-1)^{\ell(w)}\, e^{w(\rho) - \rho}
\]
equals the bar Euler product: the product over positive roots of
$(1 - e^{-\alpha})^{\mathrm{mult}(\alpha)}$ is the graded Euler
characteristic of the bar complex $B(A_X)$. For $K3 \times E$:
the denominator is the Igusa cusp form
$\Phi_{10} = \Delta_5^2$, and the root multiplicities are the
coefficients of the elliptic genus $\phi_{0,1}(\tau, z)$.

The identification operates at three levels:
(i)~\emph{genus-$1$}: the shadow Eisenstein theorem
$D_2(\cA, s) = -24\kappa_{\mathrm{ch}} \cdot \zeta(s)\zeta(s-1)$
(Volume~I) gives the genus-$1$ DT partition function via the
Borcherds lift;
(ii)~\emph{virtual}: the shadow coefficient $S_r$ at degree~$r$
matches the virtual Euler characteristic of the moduli of
$r$-dimensional sheaves;
(iii)~\emph{motivic}: the full chiral quantum group structure
($R$-matrix, coproduct, RTT) should match the motivic DT
invariants of Kontsevich--Soibelman.

This three-level correspondence is the deepest connection between
the algebraic-geometric programme of Volumes~I--II and the
enumerative geometry of this volume.

\bigskip

\noindent\textbf{Convergence and divergence.}
The shadow obstruction tower and the topological string free
energy produce genus-$g$ amplitudes with opposite analytic
behaviour. The shadow tower has Bernoulli decay
$|F_g^{\mathrm{sh}}| \sim (2\pi)^{-2g}/g$: convergent, with
radius~$2\pi$, Borel entire. The topological string has
factorial growth $|F_g^{\mathrm{top}}| \sim (2g)!$: divergent,
Gevrey-$1$, requiring non-perturbative completion. The
shadow is the algebraic soul: the part computable from OPE
structure constants alone, without worldsheet instantons.
For class~$\mathbf{G}$ algebras (Heisenberg, lattice VOAs),
the shadow is the full answer. For class~$\mathbf{M}$
(Virasoro, $\cW_{1+\infty}$), the shadow is a strict subtower,
and the instanton sum is infinite. The instanton gap
$F_g^{\mathrm{top}} - F_g^{\mathrm{sh}}$ at each genus is
the CY frontier.

The $\mathbf{G}$/$\mathbf{L}$/$\mathbf{C}$/$\mathbf{M}$
classification of Volume~I (shadow depth
$r = 2, 3, 4, \infty$) acquires a CY incarnation: the shadow
depth of $\Phi(\cC)$ is determined by the homological
complexity of~$\cC$, and the four classes reappear as
conditions on the Hochschild-to-cyclic spectral sequence of
the CY category. Class~$\mathbf{G}$ corresponds to CY
categories with semisimple Hochschild cohomology
(e.g.\ $\Coh(E)$ for an elliptic curve); class~$\mathbf{M}$
to those with infinite algebraic depth (e.g.\ $D^b(K3)$,
whose chiral algebra $\Phi(K3)$ carries the full
Virasoro-type shadow tower).

\bigskip

\noindent\textbf{The elliptic curve as universal bridge.}
The same elliptic curve~$E$ appears in six distinct roles
across the three volumes:
\begin{enumerate}[label=(\roman*),nosep]
\item \emph{Base curve} (Vol~I): the factorization algebra
 lives on $\Ran(E)$; the genus-$1$ propagator is
 $d\log E(z,w)$.
\item \emph{Sewing parameter} (Vol~I, MC5 analytic
 HS-sewing lane): the modular parameter~$\tau$ is the
 Hilbert--Schmidt sewing kernel.
\item \emph{CY fiber}: the DT partition function of
 $K3 \times E$ decomposes as a sewing of K3 fiber
 contributions along~$E$.
\item \emph{Jacobi form parameter}: the Fourier coefficients
 of $\phi_{0,1}(\tau,z)$ are root multiplicities and BPS
 degeneracies.
\item \emph{Siegel modular period}: the Igusa cusp form
 $\Delta_5$ lives on
 $\Sp_4(\Z) \backslash \bH_2$, which
 parametrizes principally polarized abelian surfaces; the
 product locus of pairs of elliptic curves is
 codimension~$1$.
\item \emph{Elliptic Hall algebra carrier}: the CY$_2$
 quantum vertex chiral group is carried on $\Coh(E)$, with
 Felder's dynamical elliptic $R$-matrix.
\end{enumerate}
Each appearance brings the same structural package: a modular
parameter, an $\SL_2(\Z)$ symmetry, a $q$-expansion, and a
sewing functor converting genus-$0$ data into genus-$1$ data.
The elliptic curve is the modular clock of the theory.

\bigskip

\noindent\textbf{The $\Omega$-background and $E_1$ stabilization.}
The $E_1$ stabilization theorem (Volume~I): for CY$_d$ with
$d \ge 3$, the Schouten--Nijenhuis bracket on the Lie conformal
algebra $L_\cC$ vanishes, so the classical structure is abelian
and all noncommutative structure arises from quantization.
The physical mechanism is the $\Omega$-background: the
equivariant parameters $\epsilon_1, \epsilon_2$ of the torus
action on $\C^2 \subset Y$ break Dunn additivity
$E_2 \simeq E_1 \otimes E_1$ by freezing one factor, producing
a single $E_1$ direction. At $\epsilon_1 = 0$, $\epsilon_2 \neq 0$
(the Nekrasov--Shatashvili limit): one $E_1$ factor survives and
the resulting algebra is the Yangian $Y(\fg)$, the quantum group
associated to the gauge theory. At generic
$(\epsilon_1, \epsilon_2)$: the quantum toroidal algebra
$U_{q,t}(\widehat{\widehat{\fg}})$ with
$q = e^{R\epsilon_1}$, $t = e^{-R\epsilon_2}$, and CY
condition $qt^{-1}(qt)^{1/2} = 1$ (for CY$_3$:
$q_1 q_2 q_3 = 1$). The passage from $E_2$ to $E_1$ is not a
loss of structure but a specialization of parameters: the
$\Omega$-background IS the $E_1$ stabilization.

This mechanism connects all three volumes: Volume~I introduces
$E_1$-chiral quantum groups as abstract algebraic-geometric
objects; this volume constructs them concretely from the
$\Omega$-background of gauge theories on CY manifolds; Volume~II
realizes their derived centres as 3d gauge theories. The
$\Omega$-background is the physical face of the $E_1$
stabilization theorem.

\bigskip

\noindent\textbf{Ten research programmes from $K3 \times E$.}
The $K3 \times E$ prototype generates ten formal research
programmes, each grounded in 90+ compute engines and detailed in
Chapter~\ref{chap:toroidal-elliptic}: (A)~constructive CY gluing
via Bridgeland charts; (B)~$\kappa_{\mathrm{cat}}$ as universal
moonshine multiplier ($A_n = \kappa_{\mathrm{cat}} \cdot
\dim(\rho_n)$); (C)~the second-quantization bridge
($\kappa_{\mathrm{BKM}}/\kappa_{\mathrm{ch}} = 5/3$);
(D)~Schottky shadow ($g \ge 4$ Torelli descent);
(E)~denominator = bar Euler product; (F)~quantum toroidal
$U_{q,t}(\widehat{\widehat{\fgl}}_1)$ as the two-parameter
chiral quantum group with CY condition $q_1 q_2 q_3 = 1$;
(G)~BPS/DT correspondence (three-level: genus-$1$, virtual,
motivic); (H)~$\Omega$-background as $E_1$ stabilization;
(I)~wall-crossing as MC gauge equivalence; (J)~the
DDCA--toroidal bridge. Each programme connects the abstract
algebraic-geometric framework of Volume~I to concrete
enumerative geometry.

\bigskip

\noindent\textbf{Maturity and scope.}
This volume is the least mature of the three. The
CY-to-chiral functor $\Phi$ is proved for $d = 2$
(Theorem~CY-A) and is a programme for $d = 3$, conditional on
the chain-level $\mathbb{S}^3$-framing. The quantum group
realization (CY-C) is conjectural: the object $C(\frakg,q)$ is
not constructed in general. The modular CY characteristic
(CY-D) is well-defined only when $A_\cC$ exists. The
$K3 \times E$ programme, with its thirteen structural results
and ten research directions, is the most developed component;
the remaining landscape (toric CY$_3$, Fukaya categories,
matrix factorizations, geometric Langlands) consists of
varying mixtures of theorems, programmes, and conjectures.
Several chapters remain stubs. What is proved is stated as
proved; what is conjectural is stated as conjecture; what is a
programme is described as such.

\bigskip

\noindent\textbf{The three volumes.}
Volume~I constructs $E_n$-chiral algebras as algebraic-geometric
objects on curves and their configuration spaces: the bar complex on
every curve geometry, $R$-matrices, Yangian structures, $E_1$-chiral
quantum groups as abstract algebraic-geometric objects, derived chiral
centres with $E_2$/$E_3$ structure, the $\SC^{\mathrm{ch,top}}$ datum
on the pair $(\cZ^{\mathrm{der}}_{\mathrm{ch}}(A),\, A)$, and the
five theorems as structural properties of the categorical logarithm.
Volume~II interprets these constructions physically: the derived chiral
centre IS the bulk of a real 3d holomorphic-topological gauge theory;
the $\SC^{\mathrm{ch,top}}$ pair IS a boundary-bulk system; 3d quantum
gravity is the climax.
This volume constructs the concrete examples: the functor
$\Phi \colon \CY_d\text{-}\Cat \to \Etwo\text{-}\mathrm{ChirAlg}$
produces chiral algebras from Calabi--Yau categories, and the
CY quantum groups (CoHA, quantum toroidal algebras, BPS algebras) are
concrete $E_1$-chiral quantum groups instantiating Volume~I's abstract
framework. The theories studied include 6d holomorphic Chern--Simons on
CY$_3$ (the primary focus), 4d and~5d HT Chern--Simons, 6d little
string theory, and M5 brane theory.

The three volumes form a triangle:
\begin{center}
\begin{tikzcd}[row sep=large, column sep=huge]
& \text{Chiral algebra } A
 \arrow[dl, "\barB"', "\text{Vol~I}"]
 \arrow[dr, "\SC^{\mathrm{ch,top}}", "\text{Vol~II}"'] \\
\text{Modular MC element } \Theta_A
 \arrow[rr, "\mathrm{av}"']
& & \text{HT field theory}
\end{tikzcd}
\end{center}
Volume~III closes the triangle from below: the CY
category~$\cC$ sits beneath all three vertices, supplying
the chiral algebra via $\Phi$, the modular trace via the CY
Euler characteristic, and the quantum group via the CoHA.

At a higher altitude, the entire programme is factorization
(co)homology at varying geometric inputs. Volume~I computes
factorization homology on the curve~$X$: the chiral algebra
$A$ is the input, the bar complex $B(A)$ on $\Ran(X)$ is the
output, and the five theorems extract genus-by-genus
invariants. Volume~II computes factorization homology on the
slab $\C \times \R$: the bar complex $B(A)$ is an $\Eone$ dg
coalgebra whose differential encodes holomorphic data along~$\C$
and whose coproduct encodes topological data along~$\R$, and the
output is the 3d holomorphic-topological field theory. This volume computes factorization homology on
the CY category~$\cC$: the cyclic bar complex
$\mathrm{CC}_*(\cC)$ is the input to~$\Phi$, and the output
is the chiral algebra that feeds Volumes~I--II. The geometric
input changes; the machine is the same.

\bigskip

\noindent\textbf{Organisation.}
Part~\ref{part:cy-categories} establishes the categorical
foundations: CY categories, cyclic $\Ainf$-structures, the
Hochschild calculus, and the
$\Eone$/$\Etwo$/$\En$ chiral hierarchy with the framing
obstruction tower. Part~\ref{part:bridge} constructs the
bridge functor $\Phi$ from CY categories to quantum chiral
algebras and proves Theorem~CY-A for $d = 2$, with the
identification of automorphic correction and shadow
obstruction tower as the central result. The heart of
Part~\ref{part:bridge} is the Drinfeld center chapter: the
categorical mechanism that lifts $\Eone$-monoidal
representation categories to $\Etwo$-braided monoidal
categories, realizing the quantum group $R$-matrix as the
half-braiding of the center construction.
Part~\ref{part:examples} works through the standard CY
landscape, centred on the $K3 \times E$ programme with its
thirteen structural results and ten research programmes.
Part~\ref{part:connections} develops the seven faces of
$r_{\mathrm{CY}}(z)$: the bar-cobar bridge to Volume~I, the
CY holographic datum, and the modular Koszul bridge.
Part~\ref{part:frontier} connects to the geometric Langlands
programme and identifies the open frontiers.

\bigskip

\noindent
Volume~I asks what the logarithm does. Volume~II asks how it
composes. Volume~III asks where it comes from.
